\section{Methods}
\label{sec:methods}

\subsection{Force closes and their formats}
A commitment transaction carries a fingerprint: the upper byte of its locktime is \texttt{0x20} and that of its single
input's sequence is \texttt{0x80}, because the two fields encode the obscured commitment number~\cite{bolt3}. We call a
transaction a force close when it has this fingerprint and spends a 2-of-2 multisignature output or a taproot output by
its key path. Its format is \emph{zero-fee} if it is a version-3 transaction with a pay-to-anchor output,
\emph{simple taproot} if it spends a taproot output, \emph{anchor} if it has at least one anchor output, and
\emph{pre-anchor} otherwise.

\subsection{Rates and cohorts}
Counts of closes follow the size of the network, so we also report rates. A public channel is \emph{open} at a height
when its funding output is still unspent. Channels whose nodes have left the network count as open until they are
closed on-chain, which is right for a rate of on-chain closes but inflates the denominator as such channels
accumulate. The force-close rate is the number of public force closes per 1,000 open public channels per month (or
week). Following our first paper~\cite{kumar2026clientdiversity}, a \emph{mass-close event} is a node party to at least
20 public force closes on one UTC day; \emph{organic} force closes are the others. The \emph{cohort incidence} at age
$a$ is the share of the public channels funded in a year that were force-closed by age $a$, among those funded at
least $a$ before the cutoff. Open channels do not enter it, so it is not biased by abandoned channels; we draw a curve
while at least half of its cohort is observed.

\subsection{The cost of a force close}
We count every transaction a force close requires, in three currencies. \emph{Fees}: the commitment's fee; an anchor
spend by its owner in the commitment's block, a CPFP fee bump, counted in full; HTLC transactions and claims; the
broadcaster's delayed \texttt{to\_local} sweep; the \texttt{to\_remote} sweep of anchor channels, which the one-block
delay makes a separate step; and later anchor spends by their owner, which we call cleanup. A spend that batches several
inputs is split equally among them; charging each close the whole transaction gives an upper bound. Not counted: the
\texttt{to\_remote} spend of a pre-anchor channel, a plain wallet output like those of a cooperative close; third-party
anchor sweeps, which are paid from the anchors' own 330~sat; and penalties, reported apart. Each year's median is set
against the median fee of a cooperative close in the same year. \emph{Time}: the days from the close until the
broadcaster's \texttt{to\_local} can be swept, from the \texttt{to\_self\_delay} revealed in its witness, and the
balance times that wait (BTC-days). To set the wait against the fees without assuming a cost of capital, the
break-even rate is the yearly return on the locked balances at which a year's BTC-days would cost as much as that
year's force-close fees. \emph{Block space}: the virtual bytes of the same transactions.

\subsection{HTLC deadlines}
\label{sec:htlc-method}
The chain gives each HTLC's expiry: an HTLC-timeout transaction's locktime is the expiry, because its pre-signed
signature commits to it, and a received HTLC's script, revealed when it is spent, contains it~\cite{bolt3}. We measure
how many blocks after its expiry each HTLC-timeout transaction confirmed and, for HTLC-success transactions, how many
blocks before. The yardstick is the margin a forwarding node keeps, its \texttt{cltv\_expiry\_delta}: we use
\HtlcMarginTight{} blocks, the tenth percentile of the values public channels advertised in our daily gossip
snapshots from \CltvDeltaFrom{} to the cutoff, and \HtlcMarginMid, their median.
The fee condition is the median fee rate of the block at the expiry. To compare formats we use the HTLCs whose
commitment had confirmed by the expiry, so that the delay is the HTLC transaction's own. Both HTLC scripts also contain
a hash of the payment hash; until point time-locked contracts replace hashes, every hop of a payment and every part of
a multi-part payment share it, so the same hash in two force closes marks one payment.

\subsection{Lifecycle cost}
For each public channel closed at least four weeks before the cutoff, the lifecycle cost adds the channel's share of
its funding transaction's fee to the cost of its close and required spends (as above, or the fee of the cooperative
close). A funding transaction's fee is split equally among the channels it opens; its other outputs count as the
opener's change and pay nothing, because the change exists only because of the open. A splice is charged once, to the
new channel's open. Rules that also charge change outputs, by value or by size, are reported as a robustness check.

\subsection{Anchors and who sweeps them}
Each P2WSH anchor is spent by its owner in the commitment's block (a bump), by its owner later (cleanup), by a third
party (an input with an empty signature, possible from the sixteenth block), or not at all. For a same-block owner
spend we also test whether the child raised the fee rate of the package, which separates a bump from cleanup that
happened to confirm in the same block. Third-party sweeps are collected chain-wide, including anchors of commitments we
never saw, and each sweep's 330~sat per anchor is split between the fee (to miners) and the outputs (to the sweeper).
The break-even fee rate, at which the fee for one anchor input equals the anchor's value, is \BreakEven~sat/vB. To
measure concentration we link sweep transactions that spend each other's outputs or pay the same address; each
connected group is an operator. Sweeps whose outputs are all \texttt{OP\_RETURN} give everything to miners and are
counted apart. We report shares only and name no operator or address.

\subsection{Fee spikes}
\label{sec:spike-method}
To test whether fee spikes bring force closes we use an event study and a regression. A \emph{spike week} has a median
block fee rate at least \SpikeRule{} times that of the previous eight weeks and at least \SpikeMinFee~sat/vB; only the
first week of a cluster counts. We compare organic public force closes in the spike week and the week after with the
four weeks before, pooled over all spikes. The regression, by week, relates the logarithm of organic public force
closes per open channel to the logarithm of the median block fee rate, with year fixed effects and a linear trend, so
that weeks are compared within the same fee regime rather than 2023 with 2025. We fit it with the fee rate of the same
week, of one to four weeks earlier and of one week later.

We judge both against placebos that keep the force-close series as observed and break only its timing against fees.
(A) For each spike, an ordinary week of the same year is drawn at random, away from every real spike
(\PlWeekDraws{} draws). (B) Each year's fee series is cut into \PlBlockWeeks-week blocks, which are shuffled within the
year (\PlBlockDraws{} draws). (C) The fee series' deviations from its year means and trend are rotated by every shift of
at least \PlMinShift{} weeks (\PlShifts{} shifts), which keeps all of its autocorrelation. A result's percentile is the
share of draws below it. The random seed (\PlSeed) and the verdict rules were fixed before the full run: a wave counts
only if the spike ratio is above 95\% of placebo~A, and a delayed rise only if the largest coefficient of lags one to
four beats the largest in each draw at $p<0.05$ under both B and C. We do not use Newey--West intervals: under the
placebos, nominal 5\% tests rejected in \PlNwRejectBlockMin--\PlNwRejectBlockMax\% of draws of B and
\PlNwRejectCircMin--\PlNwRejectCircMax\% of draws of C.

\emph{Fee guessing.} A commitment's fee rate is compared with the median and the tenth percentile of the fee rates in
its own block. \emph{Sweep timing.} A third-party sweep's block is compared with the block sixteen blocks after the
close, when the sweep first became possible. \emph{Closing bill.} The fees and block space of every close and required
spend, public or not, plus third-party sweeps, are compared with all transaction fees and block weight in each year.

\subsection{Zero-fee counterfactual}
For every anchor force close we compare the block space of the anchor mechanism, two 43-vB anchor outputs and the
transactions that spend them, with that of a zero-fee commitment: one 13-vB pay-to-anchor output, spent in the same
package, and a child with a wallet input and change for the closes that were not already bumped. An owner's anchor
input takes 69.75~vB, a third party's 51.75~vB (measured on mainnet sweeps) and a pay-to-anchor input 41.25~vB.
